% !TeX root = ../main.tex
% =================================================================================
\chapter{Requirements}\label{chap:requirements}
% =================================================================================
% TODO: This chapter is typical for theses that build an artifact. Remove or merge it if it does
% not fit your thesis.

\todo[inline]{\textbf{Optional chapter.}
  A requirements chapter is typical for theses that design and build an artifact, such as a tool, a language, or a platform.
  It may not make sense for your thesis, e.g., for an empirical study, a literature review, or a theoretical thesis.
  In that case, remove it (in \texttt{main.tex} and the thesis structure in the introduction) or merge relevant parts into other chapters.
  Discuss the structure of your thesis with your supervisor.}

This chapter describes the requirements for our approach.
\Cref{sec:requirements:elicitation} describes how we elicited them.
\Cref{sec:requirements:functional,sec:requirements:non-functional} present the functional and non-functional requirements.

% ---------------------------------------------------------------------------------
\section{Elicitation Process}\label{sec:requirements:elicitation}
% ---------------------------------------------------------------------------------

Describe the sources of your requirements, e.g., stakeholder meetings, interviews with experts, analysis of an existing tool, or literature.
For interviews, describe the participants, the interview guide (in the appendix, \cref{chap:appendix:supplementary}), and how you analyzed the answers.

% ---------------------------------------------------------------------------------
\section{Functional Requirements}\label{sec:requirements:functional}
% ---------------------------------------------------------------------------------

Give each requirement an identifier, so that you can refer to it in the concept, implementation, and evaluation, as in \cref{tab:requirements:functional}.

\begin{table}[htbp]
  \centering
  \caption{Functional requirements (example)}
  \label{tab:requirements:functional}
  \begin{tabularx}{\textwidth}{@{}lXl@{}}
    \toprule
    ID   & Requirement & Source \\
    \midrule
    FR1  & The system shall detect modifications of a file by other programs. & Interview \\
    FR2  & The system shall show conflicting modifications side by side. & Literature \\
    \bottomrule
  \end{tabularx}
\end{table}

% ---------------------------------------------------------------------------------
\section{Non-Functional Requirements}\label{sec:requirements:non-functional}
% ---------------------------------------------------------------------------------

Describe quality requirements, such as usability, performance, scalability, or maintainability, in the same way (NFR1, NFR2, \ldots).
Make them measurable where possible, so that you can evaluate them in \cref{chap:evaluation}.
